\section{The General Boundary: What Can Be Leaked Without Disturbance?}
\label{sec:information-boundary}

The two case studies suggest a broader question. Can an attacker couple a hidden
ancilla to a victim computation, learn something, and leave the victim state
unchanged?

\subsection{No information without disturbance for nonorthogonal pure states}

\begin{theorem}[Information-disturbance boundary]
Let a unitary $U$ act on a system and an environment initialized in
$\ket{0}_E$. Suppose that for two possible pure states,
\begin{equation}
  U\ket{\psi_i}_S\ket{0}_E
  =\ket{\psi_i}_S\ket{e_i}_E,
  \qquad i\in\{0,1\}.
  \label{eq:nondisturbing-map}
\end{equation}
If $\braket{\psi_0|\psi_1}\neq 0$, then
$|\braket{e_0|e_1}|=1$. Therefore the environment states are physically
indistinguishable up to a phase and contain no usable information about $i$.
\end{theorem}

\begin{proof}
Unitarity preserves inner products:
\begin{align}
  \braket{\psi_0|\psi_1}
  &=\big(\bra{\psi_0}\bra{0}\big)
    U^\dagger U
    \big(\ket{\psi_1}\ket{0}\big)\\
  &=\braket{\psi_0|\psi_1}\braket{e_0|e_1}.
\end{align}
When $\braket{\psi_0|\psi_1}\neq 0$, division gives
$\braket{e_0|e_1}=1$ up to an irrelevant phase convention. Thus the two
environment states cannot be distinguished.
\end{proof}

This compact argument is one form of the no-information-without-disturbance
principle \cite{fuchs1996information,koashi2002operations}. It explains why an
attacker cannot universally hide information about arbitrary nonorthogonal
states while preserving them exactly.

\begin{corollary}[Orthogonal data is different]
If $\braket{\psi_0|\psi_1}=0$, the theorem imposes no equality constraint on
$\ket{e_0}$ and $\ket{e_1}$. Orthogonal labels may be copied into an ancilla
without changing the source states.
\end{corollary}

For example,
\begin{equation}
  \cnot\ket{x}_S\ket{0}_E=\ket{x}_S\ket{x}_E,
  \qquad x\in\{0,1\}.
  \label{eq:classical-copy}
\end{equation}
This is a quantum circuit implementing a classical copy on the computational
basis. It is exactly the kind of operation that can leak classical control
bits, basis labels, syndrome bits, key-dependent branches, or classical
ciphertext without changing the original basis state.

\subsection{The same gate disturbs superpositions}

Let
\begin{equation}
  \ket{\psi(\theta)}
  =\cos\frac{\theta}{2}\ket{0}
  +\sin\frac{\theta}{2}\ket{1}.
\end{equation}
After a hidden CNOT to $E$,
\begin{equation}
  \ket{\Psi}_{SE}
  =\cos\frac{\theta}{2}\ket{00}
  +\sin\frac{\theta}{2}\ket{11}.
\end{equation}
Tracing out $E$ removes the off-diagonal coherence:
\begin{equation}
  \rho_S=
  \begin{pmatrix}
    \cos^2(\theta/2)&0\\
    0&\sin^2(\theta/2)
  \end{pmatrix}.
\end{equation}
The victim-state fidelity is
\begin{equation}
  F(\theta)=\bra{\psi(\theta)}\rho_S\ket{\psi(\theta)}
  =\cos^4\frac{\theta}{2}+\sin^4\frac{\theta}{2}
  =1-\frac12\sin^2\theta.
  \label{eq:fidelity-curve}
\end{equation}
At $\theta=0$ or $\pi$, the input is a classical basis state and the attack is
invisible. At $\theta=\pi/2$, the input is $\ket{+}$ and the fidelity falls to
$1/2$.

\begin{figure}[ht]
\centering
\includegraphics[width=0.78\textwidth]{figures/classical_quantum_leakage_boundary.pdf}
\caption{The classical-quantum leakage boundary for a hidden CNOT. The
computational-basis population remains unchanged for every $\theta$, but the
$X$-basis coherence is destroyed. An observer who checks only a $Z$ histogram
can miss a large state disturbance.}
\label{fig:leakage-boundary}
\end{figure}
